\documentclass{article}
\usepackage{iclr2026_conference,times}
% Explicit Times-compatible OpenType fonts for the portable XeTeX compiler.
\usepackage{fontspec}
\setmainfont{texgyretermes-regular.otf}[BoldFont=texgyretermes-bold.otf,
  ItalicFont=texgyretermes-italic.otf,BoldItalicFont=texgyretermes-bolditalic.otf]
\setmonofont{texgyrecursor-regular.otf}
\usepackage{amsmath,amssymb,booktabs,graphicx,xcolor,hyperref,url}
\iclrfinalcopy
\title{Toward a modular protease with pH control of peptide chirality preference}
\author{Research draft}
\newcommand{\CatD}{\mathrm{CatD}}
\newcommand{\ADP}{\mathrm{ADP}}
\newcommand{\angstrom}{\textup{\AA}}
\newcommand{\Dpep}{\textsc{d}}
\newcommand{\Lpep}{\textsc{l}}
\graphicspath{{figures/}}

\begin{document}
\maketitle
\lhead{Research draft in ICLR format; not submitted}

\begin{abstract}
Switching an expressible protein between cleavage of all-\Lpep{} and all-\Dpep{}
peptides requires control of stereospecific catalysis as well as suppression of
the inactive function. We investigate a modular strategy that combines short
recombinant human cathepsin D with alkaline D-peptidase in a canonical
L-amino-acid polypeptide. A preliminary pH 5.0--7.5 interval is chosen from
component-level evidence to avoid using global unfolding as the switching
mechanism. Twenty-four constructs vary a previously characterized cathepsin-D
substitution and a flexible or helical spacer. We define a computational
screen against experimental structures and an analytical requirement for
reciprocal selectivity. Independent preparation and enzymatic ligation of the
domains separates acidic maturation from final-switch characterization. A
direct peptide-fragment assay, accompanied by structural and activity-recovery
controls, can distinguish switching from irreversible damage. The first ESMFold
screen preserves both domain folds in all 24 constructs, but 22 contain severe
interdomain contacts and all have uncertain relative domain placement. Longer
flexible spacers are therefore undergoing a separately identified second screen,
alongside independent predictions and geometry checks. No fusion has been experimentally
validated, and sufficient reciprocal activity in the proposed pH interval remains
the central unresolved requirement.
\end{abstract}

\section{Introduction}
Protease stereoselectivity normally protects peptides of the opposite backbone
chirality. A useful switch would reverse which enantiomer is degraded while
leaving the other substantially intact. Here the catalyst itself must contain
only canonical L-amino acids, permitting recombinant production. This excludes
the direct construction of a mirror-image D-protease. It does not exclude a
natural L-protein whose active site recognizes D-peptides.

We consider a single final polypeptide with two catalytic domains. The task is
therefore to switch the substrate preference of the whole protein by reciprocal
control of existing activities. It is not necessary to invert the stereochemical
preference of one active site. This architecture preserves experimentally
supported catalytic machinery while concentrating uncertainty in domain
compatibility and regulation. The first target is purified in-vitro cleavage of
short aromatic peptides; general degradation of arbitrary all-D proteins is
outside the evidence available here.

Our evaluation separates three questions: whether a sequence plausibly retains
both folds, whether those folds remain intact during the input change, and
whether the resulting rates exhibit the desired opposite selectivities.
Structure prediction addresses primarily the first question. Neither a high
confidence score nor a favorable substrate docking score establishes peptide
hydrolysis or the off-target rate.

\begin{figure}[t]
\centering
\includegraphics[width=0.93\linewidth]{architecture.pdf}
\caption{Proposed architecture and input. The two catalytic functions have
separate natural precedents. Reciprocal selectivity and reversible switching
of the fusion remain hypotheses. Domain widths are schematic, not to scale.}
\label{fig:architecture}
\end{figure}

\section{Related work and scaffold selection}
\paragraph{D-peptide hydrolysis by L-proteins.}
Alkaline D-peptidase (ADP) from \emph{Bacillus cereus} DF4-B cleaves
D-phenylalanine oligomers, including an entirely D tetrapeptide, without reported
cleavage of their L counterparts in the original study \citep{asano1996}.
Its structure, 4Y7P, supports a conserved catalytic network involving Ser74,
Lys77, Tyr171, and His312 in the publication's numbering \citep{nakano2015}.
These coordinates describe an apo enzyme: the proposed peptide poses were
computed. The documented expression construct deletes precursor residues
1--25, and its activity was verified. A distinct D-aminopeptidase provides an
alternative scaffold, but its larger architecture and different preferences
are less convenient for this initial aromatic-substrate panel
\citep{bompard2000}. Other enzymes with activity on D-amino-acid amides or
mixed-chirality substrates do not automatically meet the all-D-peptide requirement
\citep{fanuel1999,haim2026}.

\paragraph{A conformationally regulated L-protease.}
Human cathepsin D has experimentally observed active and occluded conformations
\citep{baldwin1993,lee1998}. Movement of its N-terminal segment restricts access
to the active site in the structure obtained at pH 7.5. The short recombinant
form, srCatD, retains six additional N-terminal residues and can be produced
through bacterial expression, refolding, and processing \citep{beyer1996}.
Experiments with srCatD and substitutions E180Q, D187N, and Y10F connect these
positions with catalytic efficiency and the lag following a pH change
\citep{goldfarb2005}. Those measurements support the choice of variants but
do not show that they improve reciprocal selectivity in an ADP fusion.

\paragraph{Why begin with mild pH control?}
The ADP study reports stability after incubation over pH 5--10 and measurable
activity at pH 7.5, although its optimum is substantially more alkaline
\citep{asano1996}. The proposed pH 5.0--7.5 window consequently uses an overlap
in component stability rather than either enzyme's optimal rate. Stability
after incubation is distinct from instantaneous activity: the acidic ADP
leakage rate remains unknown. Likewise, cathepsin-D kinetics depend on substrate,
and a crystal grown near pH 5 is not a kinetic measurement in dilute assay buffer.
Pepsin-based switching involving very acidic activation was considered but not
advanced, given the folding requirement and the lack of supporting ADP data
there; pepsin itself also presents alkaline stability concerns \citep{konno2000}.

\paragraph{Computational design context.}
Diffusion-based backbone design and inverse folding substantially expand
accessible protein architectures \citep{watson2023,dauparas2022,dauparas2025,rfdiffusion2}.
Recent enzyme design results underscore the need to distinguish the actual
chemical reaction established experimentally from a broad enzyme-class label
\citep{lauko2025}. De novo pH-responsive assemblies and engineered protein
switches offer alternative regulatory routes \citep{boyken2019,dagliyan2018}.
For this project, redesigning an architecture around existing peptide-hydrolysis
evidence is a more conservative experimental starting point than inventing
both D-peptide catalysis and its switch simultaneously. This preference is a
design judgment, not a comparative experimental success-rate estimate.

\section{Design and quantitative requirements}
\subsection{Construct architecture}
The final ligated sequence has the form
\begin{equation*}
\mathrm{srCatD}\! -\! \mathrm{GGGGS}\! -\! \mathrm{LPET}\! -\!
\mathrm{GGG}\! -\! \mathrm{spacer}\! -\! \mathrm{GGS}\! -\! \mathrm{ADP}_{26:388}.
\end{equation*}
The srCatD domain contains FRLVTE followed by UniProt P07339 residues 65--412.
We compare WT, E180Q, D187N, and Y10F in native mature numbering. ADP uses
UniProt P94288 residues 26--388 without substitutions. Flexible spacers are
$(\mathrm{GGGGS})_n$, $n\in\{1,3,5\}$; nominally helical spacers are
$(\mathrm{EAAAK})_n$, $n\in\{2,4,6\}$. These established linker classes motivate
a geometry comparison \citep{arai2001}; the helical label does not guarantee a
rigid conformation in the final protein. The resulting 24 sequences span
737--762 residues and preserve the native catalytic residues and four CatD
disulfide pairs.

The six-residue srCatD extension is retained because its processing and kinetic
behavior have direct precedent. The L domain is placed first to leave that
N terminus free after maturation. An aromatic-free spacer reduces obvious
substrate-like motifs but cannot establish resistance to self-proteolysis.
All constructs share a common biochemical hypothesis and should not be treated
as 24 independent solutions to the catalytic problem.

\subsection{Reciprocal selectivity}
Let $a_L(p)$ and $a_D(p)$ denote initial rates per nominal mole of intact,
processed domain at pH $p$ on a matched enantiomeric substrate pair, at fixed
concentrations and terminal chemistry. These effective rates include the
state-dependent catalytically available fraction. Initially assume negligible intrinsic cross-cleavage and
independent domains. With domain ratio $r=[L]/[D]$, the two selectivities are
\begin{equation}
S_{\mathrm{low}}=\frac{r a_L(\mathrm{low})}{a_D(\mathrm{low})},\qquad
S_{\mathrm{high}}=\frac{a_D(\mathrm{high})}{r a_L(\mathrm{high})}.
\end{equation}
For both to meet a chosen engineering target $S$, the necessary interval is
\begin{equation}
\frac{S a_D(\mathrm{low})}{a_L(\mathrm{low})}\leq r\leq
\frac{a_D(\mathrm{high})}{S a_L(\mathrm{high})}.
\label{eq:ratio}
\end{equation}
For positive rates, this interval is nonempty exactly when
\begin{equation}
\frac{a_L(\mathrm{low})}{a_L(\mathrm{high})}
\frac{a_D(\mathrm{high})}{a_D(\mathrm{low})}\geq S^2.
\label{eq:range}
\end{equation}
A one-to-one fusion additionally fixes $r=1$, which must lie inside
Eq.~\ref{eq:ratio}. Increasing the concentration of the whole fusion cannot
repair an unfavorable selectivity ratio. These relations are analytical
requirements, not fitted data or predicted activities.

When cross-cleavage is non-negligible, define $a_{ij}$ for domain $i$ acting on
substrate chirality $j$. Total rates become
$v_L=[D](r a_{LL}+a_{DL})$ and
$v_D=[D](a_{DD}+r a_{LD})$. This full matrix, measured at both pH values, is
required to assess leakage. Partial maturation, unequal active fractions,
substrate competition, and coupling between domains can invalidate the
independent-domain approximation. Rates below the detection limit support
bounds, not zeros or infinite selectivity.

\begin{figure}[t]
\centering
\includegraphics[width=0.62\linewidth]{selectivity_feasibility.pdf}
\caption{Analytical feasibility under independent domains and negligible
intrinsic cross-cleavage. Blue satisfies both selectivity targets for a
one-to-one fusion; orange requires a different domain ratio; gray admits no
ratio satisfying both. The illustrative target of 100 is author-chosen.
No experimental or simulated candidate rates are plotted.}
\label{fig:feasibility}
\end{figure}

\section{Computational methods}
\subsection{References, numbering, and provenance}
Experimental structures and sequence records are downloaded with source URLs
and SHA256 hashes. CatD states use 1LYA and 1LYW; ADP uses 4Y7P.
RCSB SIFTS mappings connect each crystal polymer entity to its UniProt sequence.
The two CatD crystal chains are mapped independently: concatenating them would
incorrectly erase the segment absent after two-chain maturation. ADP author
numbering is offset by 30 from the precursor and is not an expression-boundary
annotation. Unresolved ADP residues 26--46 remain in every designed sequence.

Canonical residue composition, catalytic mappings, sequence uniqueness, and
construct boundaries are checked during generation. In srCatD, catalytic
native Asp33 and Asp231 become positions 39 and 237. ADP precursor residues
104, 107, 201, and 342 become positions 79, 82, 176, and 317 of its expressed
domain. Every inference record includes the sequence hash, model revision,
configuration, executed Git commit, SLURM job identifier, elapsed time, and
structure checksum.

\subsection{Structure predictions and diagnostic controls}
ESMFold evaluates the five isolated-domain controls and the fusion library
\citep{lin2023}. The configuration uses four recycles and one sequence
at a time. Model revision
\texttt{75a3841ee059df2bf4d56688166c8fb459ddd97a}
is staged locally for offline inference with PyTorch 2.10.0/CUDA 12.8 and
transformers 5.12.1. An independent Boltz-2 screen is being evaluated with
explicitly disabled MSA generation, three recycling steps, 200 diffusion steps,
and a recorded random seed \citep{passaro2025}. Both are single-sequence
evaluations of this artificial fusion; they do not model a pH-dependent ensemble.
The project uses typed proto-tools wrappers \citep{merchant2026}.

\subsection{Structural screen}
Each domain is aligned separately to its experimental reference using a proper
rigid Kabsch transform, with no reflection or scale change. The CatD framework
includes all resolved C-alpha atoms except native residues 1--16, the known
mobile gate. The ADP framework includes all resolved C-alpha atoms. No
score-dependent outlier removal is allowed. Complete resolved-domain deviations
on the same fit are reported separately so exclusion of the gate from fitting
does not conceal its movement.

Diagnostics include resolved-domain pLDDT, framework C-alpha RMSD, catalytic
side-chain geometry, disulfide S--S distances, interdomain heavy-atom contacts,
and partner proximity to catalytic atoms. Both experimental CatD states are
transferred into the predicted fusion frame to test geometric interference
with ADP. Such a transfer assesses compatibility of a static arrangement; it
does not predict the probability or kinetics of either state.

Prospective triage targets are domain mean pLDDT of at least 80/100, framework
RMSD approximately at most 2.5~\angstrom, and at most 0.75~\angstrom{} worsening
relative to the corresponding isolated control. Severe new overlaps and
catalytic-pocket occupation require review. Linker uncertainty is reported
separately rather than averaged into a single pass/fail score. These targets
are author-chosen diagnostics, not calibrated probabilities of biochemical
success. Discrepancies between methods and side-chain outliers require direct
inspection rather than suppression by an aggregate score.

\subsection{Computational resource controls}
Development and documentation occur in the Wynton repository; computational
work runs through CoreHPC SLURM. Network dependencies and weights are staged
on its login node because compute nodes have no internet. A serialized
submission guard counts pending as well as running campaign GPU allocations
and refuses any submission that would exceed two GPUs. CPU analyses use the
CPU partition. Source revisions are synchronized through GitHub; results,
environments, caches, and logs remain inside the designated project mirrors.
Isolated checkouts allow new analysis without changing source used by active
jobs; their script and runtime revisions are recorded separately.

\section{Results available at this draft stage}
\paragraph{Verified sequence library.}
The generator produced 24 unique fusion sequences and five isolated controls.
All use canonical residues and pass the prospective residue-mapping checks.
This establishes a reproducible screening library; it is not evidence of
expression, stability, or catalysis.

\paragraph{Experimental-state comparison.}
Reanalysis of 1LYA and 1LYW yielded a framework C-alpha RMSD of
1.918~\angstrom{} and a maximum observed gate C-alpha displacement of
29.987~\angstrom. Numerical checks verified rigid alignment without reflection,
contact-count interpretation, and residue indexing across the CatD chain break.
These computations recover the large local rearrangement described in the
structural literature. They were performed on experimental parent structures,
not newly designed fusions.

\begin{figure}[t]
\centering
\includegraphics[width=\linewidth]{native_states_ribbon.pdf}
\caption{Reanalysis of experimental CatD conformations. (a,b) PyMOL ribbon
renderings with a shared camera after fixed-framework alignment excluding
native residues 1--16. Colored segments show the gate; dark sticks mark
catalytic Asp side chains.
(c) Displacements of matched gate positions. This compares deposited crystal
structures and does not predict a pH response of a fusion.}
\label{fig:native}
\end{figure}

\paragraph{First structure screen.}
All five isolated-domain controls and 24 initial fusions completed ESMFold
prediction. Across the fusions, resolved-domain mean pLDDT was 88.07--89.81 for
CatD and 91.21--92.27 for ADP. Framework RMSDs were 0.704--0.749 and
0.888--0.959~\angstrom, respectively. No inverted or degenerate C-alpha
stereocenters were detected against the signed-volume convention calibrated
from experimental references. These favorable domain diagnostics did not
establish acceptable whole-fusion geometry: 22 of 24 raw models had at least
one interdomain heavy-atom separation below 2~\angstrom. Only WT GS5 and Y10F
GS5 lacked these contacts. All four GS5 variants accommodated the two
transferred CatD conformations without such severe overlaps, whereas the
other initial spacers did not.

Resolved interdomain PAE was 25.04--26.27~\angstrom, indicating substantial
uncertainty in relative placement. Consequently, clashes are warning signs
about these static predictions, not proof that the corresponding proteins
cannot fold. Conversely, a clash-free arrangement does not establish that
arrangement's population in solution. An adaptive second round adds eight
GS7/GS9 constructs, retaining the original domain variants and triage criteria.
Independent Boltz-2 predictions and restrained geometry checks are still in
progress; a final ten-design panel has not yet been selected.

\section{Proposed experimental characterization}
\subsection{Production separates maturation from switching}
Prepare the srCatD donor and ADP acceptor independently, then join them by
sortase-mediated transpeptidation \citep{popp2009,kobashigawa2009}.
The donor presents a C-terminal GGGGS--LPETG site followed by a removable tail;
the acceptor exposes GGG--spacer--GGS--ADP. The final Thr--Gly bond is an ordinary
peptide bond, and all retained residues are canonical. Confirm the free
oligoglycine N terminus after tag removal rather than assuming removal of
translation's initiating methionine.

The CatD production route should preserve the documented short-form maturation
and verify its FRLVTE N terminus. The engineered local processing junction is
specified in the original Table I as KPIEF$|$FRLVTE; mapping it onto the UniProt
proenzyme supplies a sequence-defined precursor proposal. Proposed purification
and ligation tags are additional engineering choices, and the mature design
FASTA is not a complete expression precursor.
Refolding and maturation occur before ADP is introduced. After ligation, remove
sortase, unligated domains, and released tails, and confirm product identity,
purity, and each module's activity. Manufacturing yield, tag tolerance, and
self-proteolysis remain experimental risks.

\subsection{Qualify a mild, fold-preserving input window}
Use purified fusion and synthetic peptides without a fluorescent reporter
protein, coupled enzyme, or protein carrier. Begin near 25~$^\circ$C with
separately prepared buffers spanning pH 5.0--7.5 in 0.5-unit increments.
Match ionic strength and solvent, verify pH at assay temperature, and compare
overlapping buffer chemistries. Introduce the protein by dilution or buffer
exchange, avoiding concentrated acid or base directly in the protein stock.

Measure isothermal far-UV CD before and after the intended assay interval,
analytical size-exclusion behavior, soluble recovery, and fragmentation by
SDS-PAGE or intact-mass analysis. Then cycle pH 7.5$\to$5.0$\to$7.5 and the
reverse sequence. Assay each module at its own reference condition with fresh
substrate and handling-matched controls. Global CD can conceal damage to one
domain, making recovery of both activities essential. A preliminary criterion
is at least 90\% activity recovery with no detectable unfolding, aggregation,
or proteolytic loss at assay precision. This is an engineering criterion,
not a guarantee of zero risk. Conditions failing qualification are excluded
before selectivity is interpreted.

\subsection{Measure all four state--chirality combinations}
Start with free-acid D-Phe$_4$, an ADP-supported substrate, and the chemically
matched L enantiomer. Establish parent CatD activity on the exact L compound
before interpreting fusion data. Historical CatD evidence for tri- and
tetraphenylalanine methyl esters does not establish identical behavior of the
free acid \citep{keilova1968}. An optional matched ester panel must distinguish
terminal ester hydrolysis from peptide-bond cleavage. Check solubility and
mass balance to avoid calling precipitation degradation.

Quantify intact substrate and hydrolysis products by LC--MS/HPLC with standards,
using initial rates for both enantiomers at both selected pH values. The two
enantiomers have identical masses; analyze them separately or use validated
chiral separation. Subsequently test mixed substrates to detect competition.
Include no-enzyme, isolated-parent, unfused-mixture, and inhibited-module or
properly processed catalytic-knockout controls. A CatD knockout cannot be
assumed to undergo the active enzyme's autocatalytic maturation. Quench only
withdrawn analytical aliquots; quenching is not the
switching input. Report time courses, switching lag, detection limits, and
rate-ratio uncertainty across independent preparations.

\section{Discussion and limitations}
The architecture retains direct evidence for catalysis of each chirality while
making a substantial new claim only as a testable hypothesis: both functions
can be gated reciprocally within the same fold-preserving interval. The most
informative early experiment is the parent-domain activity matrix in that
interval. If it fails Eq.~\ref{eq:ratio} at a one-to-one ratio, making many
linker variants is unlikely to solve the underlying rate imbalance. A variant
with greater acidic L activity can also increase high-pH leakage; greater
activity alone is not a successful switch.

Computational comparisons with native structures are useful for detecting
gross architectural damage. They cannot certify soluble expression,
disulfide formation, ligation efficiency, catalytic competence, or a
pH-dependent conformational population. Related predictors are also not
statistically independent experiments and may share structural training data.
The proposed panel is a controlled family of redesigns, not ten distinct
mechanisms. Experimental validation is needed before claiming general
chirality-switchable proteolysis or publication-level biochemical success.

The analytical framework and explicit evidence separation are the principal
methodological contributions of this draft. ICLR formatting does not imply
that this work has been submitted, accepted, or meets that venue's novelty
standard. The final report will preserve unsuccessful evaluations and identify
remaining uncertainty alongside any selected experimental panel.

\bibliography{references}
\bibliographystyle{iclr2026_conference}

\appendix
\section{Reproducibility and artifact map}
The repository's \texttt{protease\_design} directory contains the literature
review, prospective design specification, running manifest, characterization
plan, and executable scripts. Bulk reference records, inputs, predicted
coordinates, confidence data, environment snapshots, and job results are
retained under ignored project data. The manifest records actual job states
and revisions. Final sequence and ranking artifacts will be added only after
the corresponding computations have completed.

\section{AI assistance disclosure}
An OpenAI Codex assistant conducted literature searches, proposed the
architecture and screening criteria, wrote analysis and orchestration code,
and drafted this manuscript under the human request and constraints recorded
in the project. The assistant is not an author. Human researchers must review
the source interpretations, code, biological assumptions, and final claims
before experimental use or submission. No experimental data were supplied or
generated by the assistant. Any subsequently added measured data must be
identified separately from predictions and proposals.

\end{document}
